\setcounter{neje}{61}
\begin{nejer}[Repetido]
Sea $(X, \mathcal{A})$ un espacio medible y sea $\mu_n\colon \mathcal{A} \to [0, \infty]$ una sucesión de medidas. Muestre que la fórmula
$$
\mu(A) \coloneq \sum_{n=1}^{\infty} \mu_n(A)
$$
para $A \en \mathcal{A}$ define una medida $\mu\colon \mathcal{A} \to [0, \infty]$. Sea $f\colon X \to \mathbb{R}$ una función medible. Muestre que $f$ es $\mu$-integrable si y solo si
$$
\sum_{n=1}^{\infty} \int_X |f| \dd{\mu_n} < \infty.
$$
Si $f$ es $\mu$-integrable, pruebe que
$$
\int_X f \dd{\mu} = \sum_{n=1}^{\infty} \int_X f \dd{\mu_n}.
$$
\end{nejer}
\begin{sol}
\emph{Veamos que $\mu$ es medida}:\\
Por un lado, como cada $\mu_n$ es medida, $\mu_n(\varnothing) = 0$, por lo cual, $\mu(\varnothing) = \sum\mu_n(\varnothing) = 0$. Por otro lado, sea $\set{A_k}_{k\en\mathbb{N}}$ una familia contable de conjuntos de $\mathcal{A}$ disjuntos a pares. Entonces, aplicando la $\sigma$-aditividad de cada $\mu_n$:
\[
\mu\biggl( \bigcup^\infty_{k=1} A_k\biggr)
= \sum^\infty_{n=1}\mu_n\biggl( \bigcup^\infty_{k=1} A_k\biggr)
= \sum^\infty_{n=1}\sum^\infty_{k=1} \mu_n(A_k)
\]
Sin embargo, en el \myref{nejer:1.39} ya vimos que podemos conmutar estas dos series (Fubini para series), así
\[
\mu\biggl( \bigcup^\infty_{k=1} A_k\biggr)
= \sum^\infty_{n=1}\sum^\infty_{k=1} \mu_n(A_k)
= \sum^\infty_{k=1}\sum^\infty_{n=1} \mu_n(A_k)
= \sum^\infty_{k=1} \mu(A_k)
\]
\emph{Veamos que si $g\colon X \to [0,\infty]$ es medible, entonces $\int_X g\dd{\mu} = \sum_n\int_X g \dd{\mu_n}$}:\\
Si $s = \sum^{m}_{k=1} \alpha_k \chi_{A_k}$ es cualquier función escalonada medible (no negativa), entonces
\[
\int_X s \dd{\mu}
= \sum^m_{k=1} \alpha_k\mu(A_k) 
= \sum^m_{k=1} \alpha_k \sum^\infty_{n=1} \mu_n(A_k)
= \sum^m_{k=1} \sum^\infty_{n=1} \alpha_k\mu_n(A_k)
\] 
La linealidad del límite garantiza que $\sum^m_{k=1}\sum^\infty_{n=1} \alpha_k\mu_n(A_k) = \sum^\infty_{n=1}\sum^m_{k=1}\alpha_k\mu_n(A_k)$, así que
\[
\int_X s \dd{\mu}
= \sum^\infty_{n=1}\sum^m_{k=1}\alpha_k\mu_n(A_k)
= \sum^\infty_{n=1} \int_X s \dd{\mu_n}
\]
Ahora bien, si $g \colon X \to [0,\infty]$ es cualquier función medible, sabemos que existe una sucesión monótona $s_k$ de funciones escalonadas, medibles y no negativas, tal que $\lim_{n\to\infty}s_k = g$. El teorema de convergencia monótona garantizará que
\begin{equation}\label{eq:convmon}
\int_X g \dd{\mu} 
= \lim_{k\to\infty} \int_X s_k \dd{\mu}
\quad\text{y}\quad
\forall n\en\mathbb{N},\;
\int_X g \dd{\mu_n} 
= \lim_{k\to\infty} \int_X s_k \dd{\mu_n}
\end{equation}
pero ya probamos que $\int_X s_k \dd{\mu} = \sum^\infty_{n=1}\int_X s_k \dd{\mu}$; por lo tanto (\ref{eq:convmon}) implica
\begin{equation}\label{eq:limsum}
    \int_X g \dd{\mu} 
    = \lim_{k\to\infty}\sum^\infty_{n=1}\int_X s_k \dd{\mu_n}
\end{equation}
Podemos considerar, para cada $k\en\mathbb{N}$, la función auxiliar $a_k \colon \mathbb{N} \to [0,\infty]$ con $a_k(n) = \int_X s_k \dd{\mu_n}$. Notemos que, gracias a que $s_{k+1} \ge s_k$, se cumple $a_{k+1}(n) = \int_X s_{k+1} \dd{\mu_n} \ge \int_X s_{k} \dd{\mu_n} = a_k(n)$, por lo que $a_k$ es una sucesión monótona de funciones medibles (no negativas) en el espacio de medida de conteo $(\mathbb{N},2^{\mathbb{N}},\mu_c)$. Se sigue, por teorema de convergencia monótona y el hecho de que la integral  resp. a la medida de conteo coincide con la serie (\myref{nejer:1.61}), que
\[
\lim_{k\to\infty} \sum^\infty_{n=1} a_{k}(n)
= \lim_{k\to\infty} \int_\mathbb{N} a_{k} \dd{\mu_c}
= \int_{\mathbb{N}} \lim_{k\to\infty} a_k \dd{\mu_c}
= \sum^\infty_{n=1} \lim_{k\to\infty} a_k(n)
\]
es decir, $\lim_{k\to\infty}\sum^\infty_{n=1}\int_X s_k \dd{\mu_n} = \sum^\infty_{n=1}\lim_{k\to\infty}\int_X s_k \dd{\mu_n}$. Por lo tanto, (\ref{eq:limsum}) se convierte en
\[
\int_X g \dd{\mu}
= \sum^\infty_{n=1}\lim_{k\to\infty}\int_X s_k \dd{\mu_n}
\]
Como ya vimos en (\ref{eq:convmon}), $\lim_{k\to\infty}\int_X s_k \dd{\mu_n} = \int_X g \dd{\mu_n}$, entonces hemos probado que
\[
\int_X g \dd{\mu}
= \sum^\infty_{n=1} \int_X g \dd{\mu_n}
\]
\emph{Probemos lo último solicitado}:\\
Si $f\colon X \to \mathbb{R}$ es medible, por definición (y lo recién probado) $f$ es $\mu$-integrable si y solo si 
\[
\sum^\infty_{n=1} \int_X \abs{f} \dd{\mu_n}=\int_X \abs{f} \dd{\mu} < \infty
\]
Además, gracias a que $f^+$ y $f^-$ son no negativas:
\begin{align*}
\int_X f \dd{\mu} 
&= \int_X f^+ \dd{\mu} - \int_X f^- \dd{\mu}\\
&= \sum^\infty_{n=1} \int_X f^+ \dd{\mu_n} - \sum^\infty_{n=1}\int_X f^- \dd{\mu_n}\\
&= \sum^\infty_{n=1} \biggl( \int_X f^+ \dd{\mu_n} - \int_X f^- \dd{\mu_n} \biggr)\\
&= \sum^\infty_{n=1} \int_X f \dd{\mu_n}
\end{align*}
\end{sol}